\section{素数骨架：从 cusp 评价到 Hecke 刚性检验}

\subsection{cusp 处的 \texorpdfstring{$q$}{q}-展开与有限截断可审计性}

将扫描数据映射到模参数 $\tau$ 并在 cusp 区域实施 $q=e^{2\pi i\tau}$ 展开，可得到 Eisenstein 系列、$\Delta(\tau)$ 与 $j(\tau)$ 等标准对象的系数数据。由于 cusp 附近数值敏感性放大，任何有限截断必须伴随显式尾界与稳定性估计，保证“可审计”的误差控制。

\subsection{Hecke 算子与素数生成的递推约束}

Hecke 代数由素数处的算子 $T_p$ 生成；对特征形式，其 Fourier 系数满足由素数递推生成的强约束。有限数据层面可用 Sturm 型界限将“有限阶一致性”提升为全体一致性：当系数满足一定范围内的 Hecke 检验，即可将偏离视为噪声/协议失配并量化其残差。这一“素数骨架”机制为端到端协议提供算术刚性校验层，使扫描—读出的结果不仅在统计上可控，也在算术结构上可检验。

\subsection{端到端构成链条}

将前述要素组合，可形成如下构成链条（均在可操作定义内闭合）：
\[
\begin{aligned}
&\text{扫描轨道}\ \to\ \text{有限分辨率窗口读出}\ \to\ \text{差异度证书}\ \to\ \text{cusp 评价与 $q$-系数}\\
&\to\ \text{Hecke/素数骨架检验}\ \to\ \text{Zeckendorf/Ostrowski tick 编码与折叠稳定扇区}.
\end{aligned}
\]
该链条提供一种跨越动力系统、数值分析与算术约束的统一协议观：每一步均有明确输入输出与误差界，且可通过对照实验隔离其贡献。

